\subsection{I-40. A determinant-weight budget for translation-invariant matrix words}
\label{I-40}

\paragraph{Statement.}
Let $E$ be any field and let $0\ne f\in E[X_1,\ldots,X_n]$ be
homogeneous and multilinear of degree $D$, where $n\ge1$.
Assume the formal common-translation identity
\[
 f(X_1+s,\ldots,X_n+s)=f(X_1,\ldots,X_n).
 \tag{I.119}
\]
Suppose that the coefficient-cut ranks of $f$ in some order are
at most $w$. Put $m=n-2D$. Then $m\ge0$.
There is a minimal affine univariate matrix representation of
$f$ in that order, with bond sizes equal to its coefficient-cut
ranks, having the following property: every consecutive run of
generic-rank-$w$ factors in its consistent $w\times w$ padding
has length at most $m$. The same bound survives absorption of
constant factors. This assertion is valid in every characteristic;
the translation identity is a polynomial identity, not merely an
identity on the rational points of a finite field.

For $w\le3$ and $m\le1$, I-37 and I-38 consequently give
\[
 f\bigl(t\mathbf1+\mathrm{SV}_2(\mathbf y,\mathbf z)\bigr)\ne0,
 \qquad
 f\left(\left(tc_i+\frac{z_1}{a_i-y_1}
                         +\frac{z_2}{a_i-y_2}\right)_{i=1}^n\right)\ne0
 \tag{I.120}
\]
for arbitrary distinct nodes $a_i$ and arbitrary nonzero slopes
$c_i$ in the working field. No order is supplied to either map.
For this corollary use padding to width three even if the actual
width is smaller; if no rank-three core occurs, I-37 applies.

Together with P-17, the result also proves that balanced rational
two-pair evaluation hits every nonzero multilinear polynomial
of ROABP width at most three on at most nine variables, in every
characteristic and with unknown order. This last assertion is
about the balanced rational map. The proof does not identify it
with ordinary unweighted centered $\mathrm{SV}_2$ for every
nine-variable input.

\paragraph{1. The established highest-weight input and conventions.}
Write $V=E e_0\oplus E e_1$ and use the standard tensor action of
$G=\mathrm{SL}_2$. Let
\[
 u_+(s)=\begin{pmatrix}1&s\\0&1\end{pmatrix},
 \qquad h(t)=\operatorname{diag}(t,t^{-1}).
\]
We use the standard characteristic-free fact that a nonzero
$u_+$-fixed vector in a rational $G$-module, on which $h(t)$ acts
as $t^d$, has $d\ge0$. Jantzen (n.d.), Lecture~I, p.~2
(Frobenius reciprocity) and p.~3 (the explicit
$\mathrm{SL}_2$ induction computation), supply this exact input
in every characteristic; see I-B09. Upper- versus lower-Borel
conventions are changed by dualizing and conjugating by the Weyl
element. The split-torus weight decomposition is the established
integer-character grading of G-B02. These facts concern rational
group representations, and do not invoke complete reducibility
or Lie algebra eigenvalues. The convention check and an
elementary Bruhat verification are retained in the separate
\texttt{I40\_highest\_weight\_audit.md}.

Multihomogenize $f$ as the tensor
\[
 F=\sum_{S\subseteq[n]}c_S
       \bigotimes_{i=1}^n e_{\mathbf1_{i\in S}}
   \ \equiv\ \left(\prod_i U_i\right)
                f(V_1/U_1,\ldots,V_n/U_n).
 \tag{I.121}
\]
The upper unipotent fixes $e_0$ and sends $e_1$ to $e_1+s e_0$.
Thus (I.119) says exactly that $u_+(s)F=F$. Homogeneity gives
$h(t)F=t^{n-2D}F=t^mF$. The cited highest-weight fact gives
$m\ge0$. Scalar extension to an algebraic closure is harmless
here: it preserves polynomial identities and every coefficient
rank. The support spaces and matrix cores below may all be
defined over the original field $E$.

\paragraph{2. Both sides of a coefficient cut bound its determinant weight.}
Let $L_i\subseteq V^{\otimes i}$ and
$R_i\subseteq V^{\otimes(n-i)}$ be the left and right support
spaces of $F$ at the cut after position $i$. Concretely, $L_i$
is the image of contraction with all linear functionals on the
right tensor factors, and $R_i$ is defined symmetrically.
Their common dimension $r_i$ is the coefficient-cut rank;
$r_0=r_n=1$. Since $F$ is a semi-invariant of the upper Borel,
both support spaces are stable under that Borel. Indeed applying
a Borel element to a contraction is, up to its nonzero scalar
on $F$, another contraction after applying the inverse action
to the functional. Applying the inverse element gives equality
of the resulting spaces.

Let the determinant characters of the torus on these spaces be
$t^{\delta_i}$ and $t^{\delta_i'}$, respectively. They are integer
characters. The upper unipotent has determinant one on either
space: its representation matrix $G(s)$ has entries in $E[s]$
and inverse $G(-s)$, so its determinant is a unit in $E[s]$;
its value at zero is one. Consequently the determinant line of
$L_i$ is an upper-unipotent-fixed line in the rational
$G$-module $\bigwedge^{r_i}(V^{\otimes i})$. The highest-weight
fact gives $\delta_i\ge0$. The same argument on the right gives
$\delta_i'\ge0$.

In bases of $L_i$ and $R_i$, the tensor $F$ has a nonsingular
$r_i\times r_i$ coefficient matrix $C$. Torus semi-invariance
therefore gives
\[
 G_L(t) C G_R(t)^{\mathsf T}=t^m C.
\]
Taking determinants yields the exact budget
\[
 \delta_i+\delta_i'=r_i m,
 \qquad 0\le\delta_i\le r_i m.
 \tag{I.122}
\]
This is a statement in the character lattice $\mathbb Z$.
In positive characteristic, neither the inequality nor the
equality is reduced modulo the characteristic.

\paragraph{3. Minimal affine cores and their exact covariance.}
There is a natural inclusion
\[
 \iota_i:L_i\longrightarrow L_{i-1}\otimes V.
\]
To verify it, write $F$ as an element of
$L_{i-1}\otimes V\otimes V^{\otimes(n-i)}$ and contract its last
$n-i$ factors. These inclusions are Borel-equivariant. Choose
bases of the $L_i$, with the basis of $L_0=E$ equal to $1$ and
that of $L_n=EF$ equal to $F$. Write
\[
 \iota_i(b_{i,b})
   =\sum_a b_{i-1,a}\otimes
       \bigl(A_i^0[a,b]e_0+A_i^1[a,b]e_1\bigr),
 \qquad A_i(X_i)=A_i^0+X_iA_i^1.
 \tag{I.123}
\]
Repeatedly applying the inclusions reconstructs every tensor
coefficient, so the scalar matrix product
$A_1(X_1)\cdots A_n(X_n)$ equals $f$. The matrix shape is
$r_{i-1}\times r_i$, and these sizes are minimal because they
are the actual coefficient-cut ranks. This construction uses
only support-space inclusions and linear algebra, over any
field.

Let $G_i(s)$ be the upper-unipotent action matrix on $L_i$.
Equivariance of (I.123), comparing the coefficients of $e_0$
and $e_1$, gives
\[
 A_i(X_i)G_i(s)=G_{i-1}(s)A_i(X_i+s).
 \tag{I.124}
\]
In particular this formula includes both endpoints, with
$G_0=G_n=1$. It is an identity of matrices over $E[X_i,s]$;
there are no coefficient-dependent rational denominators.

Suppose $r_{i-1}=r_i=r$ and $A_i$ has generic rank $r$.
The determinant-one property of $G_{i-1}(s)$ and $G_i(s)$
implies
\[
 \det A_i(X_i+s)=\det A_i(X_i).
\]
Setting $X_i=0$ proves that this nonzero determinant is a
nonzero constant. Thus $A_i^0$ is invertible. For the torus,
comparison of the $e_0$ coefficient of (I.123) gives
\[
 A_i^0G_i(t)=t\,G_{i-1}(t)A_i^0.
\]
Taking determinants and cancelling $\det A_i^0\ne0$ proves
\[
 \delta_i=\delta_{i-1}+r.
 \tag{I.125}
\]
This is the key local restriction: a square full-rank core
increases the determinant character by its entire bond size.

\paragraph{4. The run bound and the width-three corollary.}
Pad each core consistently to $w\times w$, embedding each
$E^{r_i}$ as the first $r_i$ coordinates of $E^w$. Use the first
coordinate row and column as the scalar endpoints. This
preserves the product exactly. A padded factor can have generic
rank $w$ only if $r_{i-1}=r_i=w$ and its unpadded core is full
rank. On a run of $\ell$ such factors, (I.125) telescopes to
\[
 w\ell=\delta_{\mathrm{end}}-\delta_{\mathrm{start}}
           \le wm,
 \qquad \ell\le m,
 \tag{I.126}
\]
by (I.122) at the two bounding cuts. This proves the asserted
budget.

To absorb constant factors, multiply each intervening constant
block into the following variable-dependent factor, and absorb
the terminal block into the final column. Ranks cannot increase.
If two or more remaining full-rank factors form a consecutive
run, every constant factor removed between them must itself
have had full rank; otherwise its product would make the later
factor singular. Thus those remaining factors belong to one
full-rank run before absorption. The length of the remaining
run is no larger. Initial constant blocks and the final boundary
cannot create an extra adjacent variable factor. The same run
bound therefore holds after absorption.

For width at most three, use padding to three. If $m=0$, every
factor is singular, so I-37 applies. If $m=1$, no adjacent
variable-dependent factors have generic rank three after
absorption, so I-38 applies. This proves (I.120), with its exact
node and slope quantifiers. Its ordinary polynomial sampler
has the already proved size
\[
 (D+1)\bigl((n-1)D+1\bigr)^2\binom{D+2}{2}
 \tag{I.127}
\]
over a working field with at least
$\max\{n,D+1,(n-1)D+1\}$ elements. For $D=0$ one evaluation
suffices. A proper extension uses the explicit extension-input
convention of I-21 and I-37; the proof does not assume access to
extension evaluations when the black box only accepts base-field
points.

\paragraph{5. The nine-variable balanced two-pair theorem.}
Fix distinct nodes and let $\mathcal I_2$ denote the vanishing
ideal of balanced rational evaluation with numerator and
denominator degrees at most two, as in P-17. P-17 has already
proved the entire multilinear case on at most eight variables
and the homogeneous cubic case on at most ten variables, in
every order and characteristic. We use that stable result
without changing its statement.

Suppose that a nonzero multilinear $f\in\mathcal I_2$ on nine
variables has width at most three in some order. The established
homogeneous-degree decomposition of this ideal and the
lowest-component closure for ROABP width (P-B20, Theorem~9.13
and Proposition~8.3) give a nonzero homogeneous annihilator of
width at most three, of degree between three and six. The cubic
case of P-17 excludes degree three. Multilinear reversal
preserves both the balanced ideal and all coefficient-cut ranks,
so it excludes degree six and reduces degree five to degree
four. We may therefore suppose that $f$ is a homogeneous
quartic.

Balanced rational evaluation is closed under adding a constant.
Thus every coefficient of $f(X+s\mathbf1)$ in $s$ lies in
$\mathcal I_2$. Each coefficient of $s^j$ has degree $4-j$,
and this ideal has no nonzero multilinear component of degree
at most two. Hence
\[
 f(X+s\mathbf1)=f(X)+s g(X),
 \qquad g\in\mathcal I_2\text{ homogeneous cubic}.
 \tag{I.128}
\]
If $g\ne0$, it has the same cut-rank upper bound as $f$.
Indeed common translation is an invertible product of local
linear changes on the coefficient tensor, preserving cut ranks
over $E(s)$. With $q=s^{-1}$, the coefficient matrix of
$s^{-1}f(X+s\mathbf1)=qf(X)+g(X)$ is regular at $q=0$.
Every four-by-four minor vanishes before specialization and
therefore at $q=0$. Thus every coefficient-cut matrix of $g$
has rank at most three. P-17's nine-variable cubic theorem
forces $g=0$.

We have proved the translation identity (I.119) for this
quartic. Its weight is $m=9-2\cdot4=1$. Part~4 therefore gives
nonzero evaluation by the rational map with all slopes $c_i=1$.
That map is balanced of degrees at most two, contradicting
$f\in\mathcal I_2$. This proves the full nine-variable result.
All coefficient extractions and rank specializations in this
argument are valid in every characteristic.

\paragraph{6. An explicit finite balanced sampler.}
For any multilinear polynomial $f$ of total degree at most $D$,
write
\[
 R_i=t+\frac{z_1}{a_i-y_1}+\frac{z_2}{a_i-y_2},
 \qquad
 Q(\mathbf y)=\prod_{i=1}^n(a_i-y_1)(a_i-y_2).
 \tag{I.129}
\]
The rational map $\mathbf R$ has the same vanishing ideal as
balanced degree-two rational evaluation. One way to check the
scope is to pass to an algebraic closure: a generic quadratic
denominator has two distinct roots, and partial fractions give
exactly a constant plus two simple poles. The denominators with
distinct roots form a dense open subset in every characteristic.
Thus the parametrizations have the same formal polynomial
identities; no assertion about every rational function splitting
over the original finite field is needed.

For a nonzero pullback, $P=Qf(\mathbf R)$ is a nonzero polynomial.
Multilinearity is used here: expanding one monomial of $f$,
each index contributes either $(a_i-y_1)(a_i-y_2)$ or
\[
 t(a_i-y_1)(a_i-y_2)+z_1(a_i-y_2)+z_2(a_i-y_1).
\]
Consequently $P$ has degree at most $n$ in each selector $y_h$
and total degree at most $D$ in $(t,z_1,z_2)$. Choose $n+1$
distinct selector values, all outside $\{a_1,\ldots,a_n\}$,
and $D+1$ distinct amplitude values $b_0,\ldots,b_D$.
Sample the three amplitudes at the Newton simplex
\[
 (t,z_1,z_2)=(b_j,b_k,b_\ell),\qquad
               j,k,\ell\ge0,\quad j+k+\ell\le D.
\]
Univariate interpolation in the two selectors and the lower-set
Newton interpolation in the three amplitudes detect every
nonzero polynomial of these degrees. Every selected denominator
is nonzero. Thus the nine-variable theorem has a finite hitting
set of at most
\[
 (n+1)^2\binom{D+3}{3}
 \tag{I.130}
\]
points, for $n\le9$, over a working field with at least
$\max\{2n+1,D+1\}$ elements. The nodes and the sets of sample
values are fixed in advance and are independent of the hidden
order and all input coefficients. Extension evaluations, if
needed, obey the same explicitly stated black-box convention
as above.\hfill$\square$

\paragraph{Scope and the next gate.}
The determinant budget is not a polynomial-in-width PIT
construction. It controls runs of square full-width cores of a
homogeneous translation-invariant tensor; it does not constrain
every lower-rank transition or arbitrary nonhomogeneous words.
At weight two it permits two consecutive invertible cores, so
I-38 alone no longer closes the scalar question. The
translation-invariant ten-variable quartic is a concrete next
case, and the unrestricted mixed width-three problem and the
growing-width isolation transfer remain open here.

The highest-weight input is established representation theory.
The determinant-character budget, its core covariance
application, and the nine-variable deduction are proved above
as deductions in this research record, without a historical
priority claim. The standard dominance calculation and the
independent all-characteristic checks are recorded separately.
